% !TEX program = pdflatex
\documentclass[12pt]{article}
\usepackage[a4paper,margin=1in]{geometry}
\usepackage[T1]{fontenc}
\usepackage[utf8]{inputenc}
\usepackage{lmodern}
\usepackage{amsmath,amssymb,amsthm,mathtools,bm}
\usepackage{booktabs,array,tabularx}
\usepackage{enumitem}
\newcolumntype{L}[1]{>{\raggedright\arraybackslash}p{#1}}
\newcolumntype{Y}{>{\raggedright\arraybackslash}X}
\usepackage{setspace}
\usepackage[round,authoryear]{natbib}
\usepackage[colorlinks=true,linkcolor=blue,citecolor=blue,urlcolor=blue]{hyperref}

\onehalfspacing
\setlength{\parindent}{1.5em}
\setlength{\parskip}{0.12em}
\setlist[itemize]{leftmargin=2em,itemsep=0.25em,topsep=0.35em}
\setlist[enumerate]{leftmargin=2em,itemsep=0.25em,topsep=0.35em}
\emergencystretch=3em

\newtheorem{proposition}{Proposition}
\newtheorem{corollary}{Corollary}
\newtheorem{remark}{Remark}
\newtheorem{definition}{Definition}
\newtheorem{assumption}{Assumption}

\newcommand{\E}{\mathbb{E}}
\newcommand{\1}{\mathbf{1}}
\newcommand{\dd}{\mathrm{d}}
\newcommand{\calR}{\mathcal{R}}

\title{MAH, Commercialization Frictions, and Original-Drug Innovation}
\author{}
\date{}

\begin{document}
\maketitle
\thispagestyle{plain}

\begin{abstract}
\noindent
This paper develops a partial-equilibrium mechanism model of how China's Marketing Authorization Holder (MAH) system can affect original-drug innovation. Firms choose original-drug R\&D while anticipating the manufacturing and commercialization route available if a project advances. Before final clinical success and marketing approval are realized, a developer without sufficient internal manufacturing capacity must form a manufacturing plan and identify qualified production support, even though the formal entrusted-production arrangement and MAH approval are realized only later. MAH affects expected R\&D value by adding a regulated holder--producer separation route, lowering the route-specific policy and contractual friction of that retained external route, and changing the composite probability that a route-planning project becomes an observed approval or launch. The baseline closes only the qualified contract-manufacturing support market, generating an endogenous effective service cost $p_m^*(M,\eta)$ that can attenuate the reform's effect. The model does not treat MAH as a demand shock, a scientific-productivity shock, a removal of holder responsibility, or a universal reduction in all commercialization costs. It delivers comparative statics for route value, R\&D intensity, and observed realized original-drug outcomes, and it clarifies which data moments can discipline composite objects in a mechanism calibration.
\end{abstract}

\noindent\textit{Keywords:} MAH; original-drug innovation; commercialization frictions; entrusted production; route choice.\\
\noindent\textit{JEL codes:} D22; L65; O31; O38; I18.

\section{Introduction}

China's MAH system separates the right to hold marketing authorization from a strict requirement of in-house production. This institutional change is potentially important for original-drug innovation because many research-originating firms can generate valuable projects but lack complete manufacturing capacity. The central mechanism studied here is therefore not demand expansion or scientific productivity. It is an expected-commercialization-value channel: when a firm anticipates that an advancing project can be developed with qualified manufacturing support and later commercialized through an additional retained route, current R\&D becomes more valuable.

The mechanism builds on two established ideas. First, pharmaceutical R\&D responds to expected downstream rewards \citep{Schmookler1966,AcemogluLinn2004,Finkelstein2004,DuboisDeMouzonScottMortonSeabright2015,BudishRoinWilliams2015}. Second, an innovator's ability to appropriate returns depends on complementary assets and commercialization strategy \citep{Teece1986,AroraFosfuriGambardella2001,GansStern2003}. MAH fits this second logic because it changes how an innovator may access manufacturing capability while retaining the authorization asset.

Recent evidence on China's pharmaceutical reforms helps separate this route mechanism from two nearby channels. \citet{JiaMaYangZhang2023} study the 2015 drug-review reform and emphasize approval delay, review capacity, and the composition of firms and products entering the pipeline. \citet{BarwickXiaXia2026} study the National Reimbursement Drug List and emphasize demand and market-size incentives. Those channels can coexist with MAH, but they are not the MAH mechanism in this paper. The baseline holds scientific productivity and the market-return environment fixed and asks whether a regulated holder--producer separation route changes expected commercialization value.

For a development-economics audience, the relevant bottleneck is a regulated complementary asset. MAH is modeled as an institutional reform that can relax a manufacturing-complement constraint for research-oriented firms, while preserving holder responsibility and allowing scarcity in qualified production support to attenuate the effect.

The model is deliberately narrow. It is not a full dynamic structural model of the pharmaceutical industry. It does not solve product-market clearing, R\&D input-market clearing, free entry, exit, or an invariant firm distribution. The baseline closes only the market directly tied to the MAH mechanism: the market for qualified contract-manufacturing support. This single closure is important because MAH-induced demand for retained entrusted production support can raise the effective support cost $p_m^*(M,\eta)$, preventing the model from mechanically implying that MAH increases innovation for all firms.

The contribution is to organize the mechanism in a way that is internally consistent with feasible data. The theory separates route-planning-stage project arrival from observed approvals or launches, keeps residual holder responsibility in the entrusted-route payoff, and treats route-use data as a proxy for realized holder--producer separation rather than as the policy treatment itself. The empirical role of the model is correspondingly modest: it maps feasible approval-side moments to composite objects such as route-realization value and innovation-realization scale, without claiming to identify every primitive friction separately.

The model is presented in two layers. The main text gives the shortest complete version needed for the paper's argument. The technical appendix gives the object-status audit, payoff accounting, Bellman derivation, CMO market regularity conditions, comparative-static derivations, and calibration boundaries. Research notes are kept outside the manuscript and are used only as a model-design control panel.

\section{Institutional Background}

The relevant institutional point is timing. Under China's drug-registration framework, a firm need not already be the final commercial holder when it starts R\&D. The economic choice is forward-looking: a firm invests today while anticipating the manufacturing and commercialization environment it will face if a project advances. In practice, production support and manufacturing readiness must be planned before final approval is observed. After obtaining a Drug Approval License, the applicant becomes the MAH \citep{NMPARegistration2022}. The model therefore treats MAH as affecting expected future route feasibility and realization rather than current scientific productivity.

The MAH reform creates a regulated holder-producer separation route. The 2016 pilot allowed eligible applicants in selected regions to become holders under the pilot arrangement \citep{StateCouncilMAHPilot2016,CFDAMAHImplementation2016}. The later legal framework consolidated the principle that an MAH may manufacture by itself or entrust production to a qualified manufacturer \citep{DrugAdministrationLaw2019}. Subsequent NMPA rules require holder-side quality responsibility, assessment of the entrusted manufacturer, quality and entrustment agreements, and ongoing supervision \citep{NMPAHolderResponsibility2022,NMPAManufacturing2022,NMPAEntrustedProduction2023}.

These rules imply two modeling restrictions. First, entrusted production is not anonymous outsourcing. It is a regulated retained route in which authorization remains with the holder while production is performed by a qualified external manufacturer. Second, MAH does not eliminate holder responsibility. The holder retains quality, safety, lifecycle, and monitoring obligations. The model captures this through a residual entrusted-route burden $\mu_i^E$ and a composite realization probability $\zeta_i^E(M,p_m)$ that summarizes whether the planned retained entrusted route becomes an observed approval, launch, or realized retained outcome.

\section{Model}

\subsection{Environment and Timing}

There are two periods within the firm problem. In period $t$, firm $i$ chooses original-drug R\&D effort $x_i$ while anticipating future route feasibility. In period $t+1$, a random set of projects reaches the route-planning and manufacturing-preparation stage. Each such project chooses a route and, if it plans to use entrusted production, demands qualified production support before final clinical, regulatory, and implementation outcomes are known. Final approval, launch, and the observed holder--producer arrangement are realized only afterward.

The institutional regime is $M\in\{0,1\}$, where $M=1$ denotes legal availability of the MAH-style retained entrusted route. Conditional on $M=1$, the continuous index $\eta\in[0,1]$ measures implementation intensity: how completely the post-MAH rules reduce entrusted-route friction and improve route-specific implementation. The convention is $\eta=0$ when $M=0$. Finite differences in $M$ describe route-set expansion; derivatives in $\eta$ describe marginal post-MAH implementation. The model never differentiates the binary indicator $M$.

Firm characteristics are summarized by
\[
    \theta_i=(a_i,k_i,h_i^I,q_i^E,\mu_i^E,S_i,\tau_i^T),
\]
where $a_i$ is effective research productivity in generating route-planning-stage projects, $k_i$ indexes internal manufacturing capability, $h_i^I$ indicates whether internal production is feasible, $q_i^E$ indicates technical and regulatory eligibility for qualified entrusted production, $\mu_i^E$ is the residual holder-side burden under entrusted production, $S_i$ is the gross non-retained transfer or out-licensing value, and $\tau_i^T$ is the transfer-route friction. The indicator $q_i^E$ does not mean that a producer match has already been secured; scarcity and matching pressure are summarized by the equilibrium support cost below. Project subscripts are suppressed in the main text.

\begin{assumption}[Baseline regularity]
\label{ass:baseline_regular}
The discount factor satisfies $\beta\in(0,1)$, R\&D cost has $\kappa>0$, internal production cost $C^I(k_i)$ is weakly decreasing in internal manufacturing capability, and route-realization probabilities lie in $[0,1]$. Conditional on $M=1$, $\tau_{i\eta}^E(\eta)\le0$, $\zeta_{i\eta}^E(\eta,p_m)\ge0$, and $\zeta_{ip}^E(\eta,p_m)\le0$. In the baseline, $M$ and $\eta$ do not directly shift $a_i$, the market-return objects, or the residual burden $\mu_i^E$.
\end{assumption}

The model is partial equilibrium except for the qualified contract-manufacturing support market. It takes the market-return environment and firm-characteristic distributions as given. The object $p_m^*(M,\eta)$ is an equilibrium per-project effective cost of a qualified support package, including service scarcity and production-readiness support, not a product-market price.

\subsection{Route Set and Route Payoffs}

Conditional on a project reaching the route-planning stage, the feasible route set is
\begin{equation}
    \calR_i(M)=
    \{A,T\}\cup \{I:h_i^I=1\}\cup \{E:M=1,\ q_i^E=1\}.
    \label{eq:routeset}
\end{equation}
Here $I$ is internal production, $E$ is retained entrusted production, $T$ is a non-retained transfer, sale, or out-licensing option, and $A$ is abandonment or delay.

Let $R_i^{event}$ denote the route-independent commercialization-event return from an internally realized original drug, let $\bar R_i^E$ denote the corresponding entrusted-route return kernel, and let $v$ be the continuation value of adding one retained product or certificate to the firm's stock. For each retained route $r\in\{I,E\}$, define
\begin{equation}
    \zeta_i^r(\eta,p_m)=s_i\chi_i^r(\eta,p_m)\in[0,1],
    \label{eq:zeta_decomposition}
\end{equation}
where $s_i$ is the common downstream clinical and regulatory success component and $\chi_i^r$ is route-specific implementation conditional on that downstream success. This factorization is accounting, not separate identification: approval-side data generally discipline only the product $\zeta_i^r$.

The route-planning payoffs are
\begin{align}
    G_i^I
    &=\zeta_i^I\left[R_i^{event}+v\right]-C^I(k_i),
    \label{eq:GI}\\
    G_i^E(\eta,p_m)
    &=\zeta_i^E(\eta,p_m)\left[\bar R_i^E+v\right]
      -p_m-\tau_i^E(\eta)-\mu_i^E,
    \label{eq:GE}\\
    G_i^T&=S_i-\tau_i^T,
    \label{eq:GT}\\
    G_i^A&=0.
    \label{eq:GA}
\end{align}
The costs $C^I(k_i)$, $p_m$, $\tau_i^E$, and $\mu_i^E$ are expected planning-stage commitments and resource costs, so they are not multiplied by $\zeta_i^r$. If a component is paid only after realization, its success-contingent expectation is included in the corresponding cost object. This convention prevents success probabilities from being counted twice. The transfer route is non-retained; abandonment or delay is the normalized outside option.

The entrusted realization object depends on technical eligibility, route friction, firm characteristics, and support scarcity. The baseline imposes
\[
    \frac{\partial\zeta_i^E}{\partial p_m}\le0,
    \qquad
    \frac{\partial G_i^E}{\partial p_m}
    =(\bar R_i^E+v)\frac{\partial\zeta_i^E}{\partial p_m}-1<0.
\]
The transfer value $S_i-\tau_i^T$ is kept exogenous in the baseline. A Grossman-Hart-Moore-style outside-option interpretation is useful, but it is not added as a main state variable \citep{GrossmanHart1986,HartMoore1990}.

\subsection{Route Choice and CMO Market Closure}

Planned route choice follows a logit model throughout. Project-route fit shocks are independent Type-I extreme value and are location-normalized so that the expected maximum is exactly the log-sum expression below:
\begin{equation}
    P_i(r\mid M,\eta,p_m)
    =
    \frac{\exp(G_i^r(M,\eta,p_m)/\sigma_r)}
    {\sum_{\ell\in\calR_i(M)}\exp(G_i^\ell(M,\eta,p_m)/\sigma_r)},
    \qquad \sigma_r>0.
    \label{eq:logit_prob}
\end{equation}
The inclusive value of a route-planning opportunity is
\begin{equation}
    \Gamma_i(M,\eta,p_m)
    =
    \sigma_r\log
    \sum_{r\in\calR_i(M)}
    \exp(G_i^r(M,\eta,p_m)/\sigma_r).
    \label{eq:Gamma}
\end{equation}
The route shocks represent project-specific implementation fit or information revealed at route planning. Their common location is fixed rather than left free because $\Gamma_i$ affects the level of R\&D. The deterministic maximum is the limiting case as $\sigma_r\rightarrow0$, not the baseline model.

At any candidate support cost $p_m$, the R\&D response induced by the inclusive value is
\begin{equation}
    x_i^*(M,\eta,p_m)=\frac{\beta a_i}{\kappa}\Gamma_i(M,\eta,p_m).
    \label{eq:xstar_offeq}
\end{equation}

Each project planning route $E$ requires one qualified support package. Total demand includes background demand from non-MAH uses, which keeps the support market active before the reform:
\begin{equation}
    D_m^{tot}(p_m;M,\eta)
    =
    D_m^B(p_m)
    +\int
    a_i x_i^*(M,\eta,p_m)
    P_i(E\mid M,\eta,p_m)
    \dd H_i .
    \label{eq:cm_demand}
\end{equation}
Equation \eqref{eq:cm_demand} is planning-stage demand. It does not multiply by $\zeta_i^E$ because support is reserved before final realization; observed retained entrusted output does multiply by $\zeta_i^E$.
Qualified producer supply is upward sloping,
\begin{equation}
    S_m(p_m)=\int s(p_m,z_j^C)\dd H_C(z_j^C),
    \qquad
    \frac{\partial s(p_m,z_j^C)}{\partial p_m}\ge0.
    \label{eq:cm_supply}
\end{equation}

\begin{assumption}[CMO market regularity]
\label{ass:cmo_regular}
The support cost lies in $[0,\bar p_m]$. Total demand and supply are finite and continuous. The excess-demand function $Z_m(p;M,\eta)=D_m^{tot}(p;M,\eta)-S_m(p)$ is strictly decreasing, with $Z_m(0;M,\eta)\ge0$ and $Z_m(\bar p_m;M,\eta)\le0$. Where derivatives are used, $D_{mp}^{tot}<S_{mp}$.
\end{assumption}

\begin{proposition}[Existence and uniqueness of the CMO support cost]
\label{prop:cmo_existence}
Under Assumption \ref{ass:cmo_regular}, there is a unique $p_m^*(M,\eta)\in[0,\bar p_m]$ satisfying
\begin{equation}
    D_m^{tot}(p_m^*;M,\eta)=S_m(p_m^*).
    \label{eq:baseline_cm_clearing}
\end{equation}
\end{proposition}

\noindent Existence follows from continuity and the boundary inequalities; strict monotonicity gives uniqueness. No other market is closed.

When empirical predictions distinguish regions, dosage forms, or technology segments, $m$ indexes a separate qualified-support market: the demand integral uses the cell-specific innovator distribution $H_m$, and supply uses $H_{C,m}$. The one-market baseline is the special case with a single $m$; cross-cell substitution is not modeled unless explicitly added.

\begin{definition}[Baseline partial equilibrium]
\label{def:partial_equilibrium}
Given $(M,\eta)$ and the exogenous payoff, technology, and distributional objects, a baseline partial equilibrium is a tuple $\{p_m^*,P_i,\Gamma_i,x_i^*,D_m^{tot},S_m\}$ satisfying the route-payoff equations, logit route choice, the R\&D first-order condition, and equation \eqref{eq:baseline_cm_clearing}.
\end{definition}

\subsection{R\&D Choice and Observed Outcomes}

R\&D effort produces projects that reach the route-planning stage according to
\begin{equation}
    \lambda_i^{plan}=a_i x_i.
    \label{eq:planning_arrival}
\end{equation}
This is not a count of already approved or clinically successful drugs. Downstream success is represented by $\zeta_i^r$. R\&D effort has cost
\begin{equation}
    C_i^R(x_i)=\frac{\kappa}{2}x_i^2,\qquad \kappa>0.
    \label{eq:rd_cost}
\end{equation}
At the market-clearing support cost, the firm solves
\[
    \max_{x_i\ge0}
    \beta a_i x_i\Gamma_i(M,\eta,p_m^*(M,\eta))-\frac{\kappa}{2}x_i^2 .
\]
The interior solution is the equilibrium evaluation of equation \eqref{eq:xstar_offeq}:
\begin{equation}
    x_i^*(M,\eta)
    =\frac{\beta a_i}{\kappa}\Gamma_i(M,\eta,p_m^*(M,\eta)).
    \label{eq:xstar}
\end{equation}
The second derivative is $-\kappa<0$, so the solution is unique whenever the non-negativity constraint does not bind. Because route $A$ has normalized value zero, $\Gamma_i\ge0$.

Let $\calR_i^{ret}(M)=\calR_i(M)\cap\{I,E\}$. Observed retained output differs from planning-stage project arrival:
\begin{equation}
    \lambda_i^{obs,ret}(M,\eta)
    =
    a_i x_i^*(M,\eta)
    \sum_{r\in\calR_i^{ret}(M)}
    P_i(r\mid M,\eta,p_m^*)\zeta_i^r(M,\eta,p_m^*).
    \label{eq:lambda_obs}
\end{equation}
The notation $p_m^*$ in downstream equations means $p_m^*(M,\eta)$. Transfer and abandonment are excluded from this retained-product outcome. A broader data outcome can be represented by explicit route-specific observation weights, but its denominator must match the sample definition.
For the retained entrusted route,
\begin{equation}
    \lambda_i^{obs,E}(M,\eta)
    =
    a_i x_i^*(M,\eta)
    P_i(E\mid M,\eta,p_m^*)\zeta_i^E(M,\eta,p_m^*).
    \label{eq:lambda_obs_E}
\end{equation}

\section{Theoretical Implications}

\begin{proposition}[Discrete route-set expansion]
\label{prop:set_expansion}
For an eligible project with $q_i^E=1$, hold $p_m$, $\eta$, and all old-route payoffs fixed. Let
\[
    Q_i^0=\sum_{r\in\calR_i(0)}\exp(G_i^r/\sigma_r).
\]
Adding route $E$ raises the inclusive value by
\begin{equation}
    \Delta\Gamma_i^{set}
    =\sigma_r\log\left[
    1+\frac{\exp(G_i^E(\eta,p_m)/\sigma_r)}{Q_i^0}
    \right]>0.
    \label{eq:set_expansion}
\end{equation}
\end{proposition}

\noindent This is a finite comparison in the binary institution $M$. It is distinct from the marginal effect of improving implementation after route $E$ has become legally available.

\begin{proposition}[Implementation intensity and CMO attenuation]
\label{prop:cmo_attenuation}
Conditional on $M=1$, suppose only route $E$ depends directly on $\eta$. At a differentiable interior CMO equilibrium,
\begin{align}
    G_{i\eta}^E
    &=(\bar R_i^E+v)\zeta_{i\eta}^E-\tau_{i\eta}^E\ge0,
    \label{eq:GE_eta}\\
    G_{ip}^E
    &=(\bar R_i^E+v)\zeta_{ip}^E-1<0,
    \label{eq:GE_p}\\
    \frac{\dd p_m^*}{\dd\eta}
    &=\frac{D_{m\eta}^{tot}}{S_{mp}-D_{mp}^{tot}}\ge0
    \quad\text{if }D_{m\eta}^{tot}\ge0,
    \label{eq:dpdeta}\\
    \frac{\dd\Gamma_i}{\dd\eta}
    &=P_i(E)\left[
      G_{i\eta}^E+G_{ip}^E\frac{\dd p_m^*}{\dd\eta}
      \right].
    \label{eq:dGammadeta}
\end{align}
Consequently, $\dd x_i^*/\dd\eta=(\beta a_i/\kappa)\dd\Gamma_i/\dd\eta$. The first term in brackets is the direct implementation gain; the second is weakly negative and is the equilibrium support-cost attenuation.
\end{proposition}

\noindent Equation \eqref{eq:dpdeta} follows by differentiating market clearing while holding the exogenous background-demand and supply schedules fixed. A positive overall innovation response requires the direct gain to exceed this price feedback; the model does not impose that result mechanically.

\begin{proposition}[Observed entrusted-route decomposition]
\label{prop:obs_decomposition}
The change in observed retained entrusted output can be decomposed as
\begin{align}
    \dd \lambda_i^{obs,E}
    =&\;
    a_i P_i(E)\zeta_i^E\dd x_i^* \notag\\
    &+
    a_i x_i^*\zeta_i^E\dd P_i(E) \notag\\
    &+
    a_i x_i^*P_i(E)\dd\zeta_i^E.
    \label{eq:obs_decomposition}
\end{align}
\end{proposition}

\noindent Observed approvals or launches may rise because firms invest more in R\&D, route choice shifts toward $E$, or route realization improves. Approval-side data do not separately identify these components without application, status, and project-level information.

\begin{proposition}[Research and manufacturing heterogeneity]
\label{prop:heterogeneity}
For firms with feasible internal production, if only $G_i^I$ depends on $k_i$, then
\begin{equation}
    \frac{\partial G_i^I}{\partial k_i}=-C_k^I(k_i)>0,
    \qquad
    \frac{\partial P_i(E)}{\partial k_i}
    =-\frac{P_i(E)P_i(I)}{\sigma_r}
      \frac{\partial G_i^I}{\partial k_i}<0.
    \label{eq:dPEdk}
\end{equation}
At fixed $p_m$, the direct R\&D response to implementation intensity is
\begin{equation}
    \left.\frac{\partial x_i^*}{\partial\eta}\right|_{p_m}
    =\frac{\beta a_i}{\kappa}P_i(E)G_{i\eta}^E.
    \label{eq:direct_x_eta}
\end{equation}
Thus high research productivity scales the response, while weak internal manufacturing raises entrusted-route relevance through a larger $P_i(E)$.
\end{proposition}

\begin{corollary}[Aggregate implication]
For a fixed firm distribution, aggregate route-planning-stage project arrivals rise when the positive $\Gamma_i(M,\eta,p_m^*)$ response among relevant firms is not offset by CMO support-cost feedback and the affected set has positive measure. Aggregate observed outcomes additionally require that route-use and realization terms do not fall enough to offset the planning-stage gain. This is a private-value and realized-output implication, not a welfare theorem.
\end{corollary}

\section{Empirical Predictions and Calibration Implications}

The model delivers three empirical predictions that are useful before calibration. First, MAH should increase holder--producer separation among realized original-drug projects primarily where the retained entrusted route is feasible and valuable. Second, realized original-drug outcomes should respond more for firms with stronger research capability but weaker internal manufacturing capability, provided they are technically eligible for qualified support. Third, the response should be attenuated where support is scarce. For that prediction the support market is indexed by cell $m$---for example, region, dosage form, or technology segment---and each cell has its own $p_m^*(M,\eta)$. These predictions distinguish the MAH route-friction channel from a generic demand shock, scientific-productivity shock, or approval-speed reform.

The model is intended to discipline a mechanism calibration, not to identify every primitive friction. The key empirical distinction is between route-planning-stage project arrival and observed realized outcomes. Approval data, holder--producer links, and original-drug counts are close to realized outcomes. They are not direct observations of project arrival, scientific success, or route-specific legal and managerial costs.

The observed entrusted share must be computed from realized retained products, not from an unweighted average of route-choice probabilities. For a comparable sample,
\begin{equation}
    s_E^{obs}(M,\eta)
    =
    \frac{\int a_i x_i^*P_i(E)\zeta_i^E\dd H_i}
    {\int a_i x_i^*\sum_{r\in\calR_i^{ret}(M)}P_i(r)\zeta_i^r\dd H_i}.
    \label{eq:observed_split_share}
\end{equation}
Only in a homogeneous firm or narrow cell with exactly two retained routes $I$ and $E$ does the observed conditional log odds reduce to
\begin{equation}
    \log\frac{s_{iE}^{obs}}{1-s_{iE}^{obs}}
    =\frac{G_i^E-G_i^I}{\sigma_r}
     +\log\frac{\zeta_i^E}{\zeta_i^I}.
    \label{eq:observed_log_odds}
\end{equation}
Because $\zeta_i^r=s_i\chi_i^r$, the common success component cancels in this individual or narrow-cell ratio, leaving $\log(\chi_i^E/\chi_i^I)$. It does not generally cancel before aggregation.
With firm heterogeneity, the log odds of the aggregate share in equation \eqref{eq:observed_split_share} is not a representative-firm payoff gap. It must be computed by integrating or simulating the aggregate moment.

\begin{table}[htbp]
\centering
\footnotesize
\caption{Model Objects and Feasible Data Moments}
\label{tab:model_data_mapping}
\begin{tabularx}{\textwidth}{L{0.28\textwidth}Y Y}
\toprule
Model object & Feasible moment & Claim boundary \\
\midrule
Equation \eqref{eq:observed_split_share} & Holder--producer split among comparable realized retained original drugs & Disciplines an aggregate route-realization composite, not a primitive $\tau_i^E$, $\mu_i^E$, or $\zeta_i^E$ separately. \\
$a_i x_i^*(M,\eta)$ and realization scale & Realized original-drug counts & Disciplines project-arrival and realization scale jointly, not planning-stage arrivals alone. \\
$p_m^*(M,\eta)$ or CMO scarcity & Qualified entrusted-producer participation or capacity proxy by support-market cell & Needed to interpret attenuation; direct support costs may be unavailable. \\
Entry-cost distribution $F_e$ & Firm-year research-oriented entry panel & Future or sensitivity module unless credible entry data are merged. \\
\bottomrule
\end{tabularx}
\end{table}

Three calibration boundaries follow. First, route-use moments should be interpreted as evidence about relative route attractiveness and realization among comparable retained products, not as a direct estimate of the MAH policy treatment. Second, original-drug counts discipline a composite scale combining R\&D effort, research productivity, route choice, and realization probabilities. Third, entry should not be a baseline calibrated outcome unless the data contain a credible firm-year entry panel and pre-reform capability measures.

The baseline empirical exercise should therefore focus on approval-side realized-product moments: realized original-drug counts, original-drug shares, holder-producer splits, and heterogeneity by pre-reform manufacturing capability where available. Application-to-approval conversion, true project-level success probabilities, CMO support-cost proxies, and research-oriented entry are natural future modules, but they should not be treated as already identified by approval-side records alone.

The formal appendix documents the corresponding moment equations, route-share log-odds mapping, and entry-module boundary. This separation is intentional: the main paper states what the feasible moments can support, while the appendix records the algebra behind that claim.

\section{Conclusion}

The paper's main mechanism is simple: MAH changes the value of R\&D by changing the manufacturing and commercialization route that an advancing project can plan to use before final realization. The retained entrusted route matters most for firms that can generate projects but lack complete internal manufacturing capability, and its value depends on route feasibility, residual holder responsibility, composite realization probability, outside options, and the equilibrium effective cost of qualified manufacturing support. The model remains intentionally partial equilibrium. This restriction is useful because it keeps the theory aligned with the paper's empirical objective: explaining how a specific institutional change can alter expected route value, R\&D incentives, and observed realized original-drug outcomes without turning the project into a full industry equilibrium or welfare model.

\clearpage
\bibliographystyle{aer}
\bibliography{mah_route_indicator_friction_refs}

\end{document}
